\title{Group Sequence Policy Optimization}

\begin{abstract}
We present Group Sequence Policy Optimization (GSPO), a reinforcement learning algorithm that offers enhanced stability, efficiency, and effectiveness for training large language models. Differing from traditional methods that utilize token-based importance ratios, GSPO leverages sequence-level likelihoods to establish its importance ratio, implementing clipping, reward assignment, and optimization directly at the sequence level. Our results show that GSPO outperforms the GRPO algorithm in terms of both sample efficiency and training results, markedly stabilizes Mixture-of-Experts (MoE) RL procedures, and has the capacity to streamline RL infrastructure architectures. These strengths of GSPO have played a key role in the notable advancements observed in the most recent Qwen3 models.
\end{abstract}

\section{Introduction}
\label{sec:intro}

Reinforcement learning (RL) has become a crucial methodology in advancing the scalability of language models \citep{o1,dpsk-r1,qwq32b,qwen3}. Leveraging extensive RL training allows language models to acquire skills necessary for handling complex tasks—ranging from advanced mathematics to high-level programming—by enabling extended and intricate chains of reasoning.

For RL to be successfully scaled alongside increasing computational resources, ensuring stable and resilient training processes is of paramount importance. Nevertheless, the prevailing RL techniques, such as GRPO \citep{grpo}, are plagued by significant instability during the training of large-scale language models. These shortcomings frequently culminate in disastrous and irreversible deterioration of model performance, as highlighted in recent works \citep{qwen3, minimax-m1}, posing a substantial barrier to further enhancing the capabilities of language models via prolonged RL training.

This study reveals that the inherent instability in GRPO originates from a fundamental misapplication and subsequent breakdown of importance sampling weights within its algorithmic structure. Such misuse results in elevated variance during training, which accumulates as the length of generated responses increases, and is further intensified by the clipping strategy—eventually leading to the failure of the model.

To overcome these principal weaknesses, we introduce \textbf{Group Sequence Policy Optimization (GSPO)}, a novel reinforcement learning method designed for training large language models. GSPO’s principal contribution is its theoretically robust importance ratio formulation, grounded in sequence likelihood \citep{click} and fully consistent with foundational importance sampling principles. Furthermore, GSPO leverages normalized rewards calculated as the advantages across multiple response samples to each query, thus preserving coherence between sequence-level reward assignment and optimization processes.

Our experimental results indicate that GSPO markedly outperforms GRPO with respect to training reliability, effectiveness, and overall outcomes. Notably, GSPO naturally addresses the instability issues that commonly arise during reinforcement learning on large Mixture-of-Experts (MoE) architectures, obviating the requirement for elaborate stabilization techniques, and offers prospects for streamlining RL frameworks. These advantages have directly enabled substantial enhancements in the recent iterations of the Qwen3 models. We anticipate that GSPO will serve as a resilient and scalable core for future progress in large-scale RL training of language models.

\section{Preliminaries}

\paragraph{Notation}
Throughout this work, we represent an autoregressive language model, governed by parameters $\theta$, as a policy $\pi_\theta$. The variable $x$ refers to a query, while the collection of all queries is labeled as $\mathcal{D}$. For any response $y$ generated in reply to a query $x$, its probability according to the policy $\pi_\theta$ is written as $\pi_\theta (y | x)=\prod_{t=1}^{|y|} \pi_\theta (y_t | x, y_{<t} )$, with $|y|$ indicating the total number of tokens in $y$. Additionally, we assign a reward $r(x, y) \in [0, 1]$ to each query-response pair $(x, y)$ by evaluating it with a verifier $r$.

\paragraph{Proximal Policy Optimization (PPO)} Using samples generated from the old policy $\pi_{\theta_\text{old}}$, PPO \citep{ppo} constrains the policy update within a proximal region of the old policy through the clipping mechanism. Specifically, PPO employs the following objective for policy optimization (we omit the KL regularization term hereinafter for brevity, as it is not the focus of this paper):

\begin{align}
\mathcal{J}_\text{PPO}(\theta) = \mathbb{E}_{ x \sim \mathcal{D},\, y \sim \pi_{\theta_\text{old}}( \cdot | x) }
\left[ \frac{1}{|y|} \sum_{t=1}^{|y|} 
\min \left( w_{t}(\theta) \widehat{A}_{t},  \, \mathrm{clip} \left( w_{t}(\theta), 1 - {\varepsilon}, 1 + {\varepsilon}\right) \widehat{A}_{t} \right)
\right],
\end{align}

The token $y_t$'s importance ratio is given by $
w_{t}(\theta) = \frac{ \pi_{\theta} (y_{t} | x, y_{<t}) }{ \pi_{\theta_\text{old}} (y_{t} | x,y_{<t})} $, where $\widehat{A}_{t}$ denotes the advantage associated with $y_t$, computed using a separate value model, and $\varepsilon$ signifies the clipping threshold for the importance ratios.

A significant limitation of PPO in real-world applications stems from its substantial dependence on the value model. In practice, the value model often matches the policy model in complexity and scale, thereby increasing both computational and memory demands. Moreover, PPO's performance is tightly linked to the accuracy of the value estimates provided. Developing an adequately precise value model is already a formidable task, but making the value model adaptable to longer outputs and more intricate scenarios exacerbates this difficulty even further.

\paragraph{Group Relative Policy Optimization (GRPO)} GRPO \citep{grpo} bypasses the need for the value model by computing the relative advantage of each response within a group of responses to the same query. Specifically, GRPO optimizes the following objective:

\begin{align}
\label{equ:grpo}
\mathcal{J}_\text{GRPO}(\theta) = \mathbb{E}_{ x \sim \mathcal{D},\, \{y_i\}_{i=1}^G \sim \pi_{\theta_\text{old}}( \cdot | x) }
\left[ \frac{1}{G} \sum_{i=1}^{G} \frac{1}{|y_i|} \sum_{t=1}^{|y_i|} 
\min \left( w_{i,t}(\theta) \widehat{A}_{i,t},  \, \mathrm{clip} \left( w_{i,t}(\theta), 1 - {\varepsilon}, 1 + {\varepsilon}\right) \widehat{A}_{i,t} \right)
\right],
\end{align}

where $G$ is the number of generated responses to each query $x$ (i.e., the group size), and the importance ratio $w_{i,t}(\theta)$ and advantage $\widehat{A}_{i,t}$ of token $y_{i,t}$ are:

\begin{align}
    w_{i,t}(\theta)=\frac{ \pi_{\theta} (y_{i,t} | x, y_{i,<t}) }{ \pi_{\theta_\text{old}} (y_{i,t} | x,y_{i,<t})},\quad\ 
    \widehat{A}_{i,t} = \widehat{A}_{i} = \frac{r(x, y_i) - \mathrm{mean} \left( \{ r(x, y_i) \}_{i=1}^G \right) }{ \mathrm{std} \left( \{ r(x, y_i) \}_{i=1}^G \right) },
\end{align}

respectively, where all the tokens in $y_i$ share the same advantage as $\widehat{A}_{i}$.

\section{Motivation}
\label{sec:motivation}

As models increase in size, incorporate greater sparsity (as in Mixture-of-Experts architectures), and generate longer responses, it becomes necessary to employ large rollout batch sizes to ensure efficient hardware usage during reinforcement learning. To enhance the efficiency of sample utilization, practitioners typically divide these extensive rollout batches into several smaller mini-batches to perform gradient updates. This practice results in an off-policy learning context, given that the sampled responses $y$ are drawn from a previous policy $\pi_{\theta_\text{old}}$ instead of the present, optimized policy $\pi_\theta$. Such a scenario underlines the importance of the clipping mechanisms used in PPO and GRPO, which are designed to limit the influence of samples that are excessively ``off-policy'' during gradient computation.

While mechanisms like clipping aim to manage this off-policy discrepancy, we identify a more fundamental issue in GRPO: \textit{its objective is ill-posed}. This problem becomes particularly acute when training large models on long-response tasks, leading to catastrophic model collapse. The ill-posed nature of the GRPO objective stems from a misapplication of importance sampling weights. The principle of importance sampling is to estimate the expectation of a function $f$ under a target distribution $\pi_\text{tar}$ by re-weighting samples drawn from a behavior distribution $\pi_\text{beh}$:

\begin{align}
\mathbb{E}_{z \sim \pi_\text{tar}} \left[ f(z) \right] = \mathbb{E}_{z \sim \pi_\text{beh}} \left[ \frac{ \pi_\text{tar} (z) }{ \pi_\text{beh} (z)} \, f(z) \right].
\end{align}

A key requirement for effective correction of the distributional gap using importance weighting $\frac{ \pi_\text{tar} (z) }{ \pi_\text{beh} (z)}$ is to average the weights across a large number of samples ($N \gg 1$) drawn from the behavior policy $\pi_\text{beh}$.

Conversely, GRPO utilizes an importance weighting factor $\frac{ \pi_{\theta} (y_{i,t} | x, y_{i,<t}) }{ \pi_{\theta_\text{old}} (y_{i,t} | x,y_{i,<t})}$ applied at every token position $t$. Because this ratio is calculated from only a single token $y_{i,t}$ sampled from the subsequent-token distribution $\pi_{\theta_\text{old}}( \cdot | x, y_{i, <t})$, it does not effectively mitigate the mismatch between distributions. Rather, it tends to inject significant variance into gradient estimates, which accumulates through lengthy sequences and becomes worse due to the application of clipping. Our empirical studies have shown this instability frequently manifests as irreversible model collapse. Should this collapse occur, continuing training—even by reverting to previous checkpoints, carefully retuning hyperparameters such as clipping thresholds, increasing sequence length, or changing RL queries—fails to restore proper model performance.

These findings highlight an inherent flaw in GRPO’s framework. Specifically, the ineffectiveness of the per-token importance weighting underscores a central concept: \textbf{the granularity of the optimization objective must align with the reward’s granularity}. Since the reward signal is assigned for the complete sequence, implementing off-policy corrections on a per-token basis is inherently misaligned. This realization compels us to abandon the token-level objective and instead investigate the application of importance weights and optimization on the entire \textit{sequence level}.

\section{Algorithm}

\subsection{GSPO: Group Sequence Policy Optimization}

Although employing the token-level importance weight $\frac{ \pi_{\theta} (y_{i,t} | x, y_{i,<t}) }{ \pi_{\theta_\text{old}} (y_{i,t} | x,y_{i,<t})}$ presents difficulties within GRPO, we note that, in the realm of language generation, the \textit{sequence-level} importance weight $\frac{ \pi_{\theta} (y | x) }{ \pi_{\theta_\text{old}} (y | x)}$ holds a robust theoretical interpretation: it quantifies the extent to which a response $y$, drawn according to $\pi_{\theta_\text{old}} (\cdot | x)$, diverges from the distribution induced by $\pi_{\theta} (\cdot | x)$. This directly corresponds with the sequence-level reward and offers an informative basis for the implementation of the clipping mechanism.

Based on this straightforward observation, we propose the \textbf{Group Sequence Policy Optimization (GSPO)} algorithm. GSPO employs the following sequence-level optimization objective:

\begin{align}
\mathcal{J}_\text{GSPO} (\theta) =
\mathbb{E}_{ x \sim \mathcal{D},\, \{y_i\}_{i=1}^G \sim \pi_{\theta_\text{old}}( \cdot | x) }
\left[ 
\frac{1}{G} \sum_{i=1}^{G}
\min \left( s_{i}(\theta)  \widehat{A}_{i},  \, \mathrm{clip} \left( s_{i}(\theta), 1 - {\varepsilon}, 1 + {\varepsilon} \right) \widehat{A}_{i} \right) 
\right],
\label{equ:gspo}
\end{align}

where we adopt the group-based advantage estimation:

\begin{align}
\widehat{A}_{i} = \frac{r(x, y_i) - \mathrm{mean} \left( \{ r(x, y_i) \}_{i=1}^G \right) }{ \mathrm{std} \left( \{ r(x, y_i) \}_{i=1}^G \right) },
\end{align}

and define the importance ratio $s_{i}(\theta)$ based on sequence likelihood \citep{click}:

\begin{align}
s_{i}(\theta) = \left( \frac{ \pi_{\theta} (y_i | x) }{ \pi_{\theta_\text{old}} (y_i | x)} \right)^{\frac{1}{|y_i|}}
=
\exp \left( \frac{1}{|y_i|} \sum_{t=1}^{|y_i|} \log \frac{ \pi_{\theta} (y_{i,t} | x, y_{i,<t}) }{ \pi_{\theta_\text{old}} (y_{i,t} | x,y_{i,<t})} \right).
\end{align}

As a result, GSPO implements clipping across complete responses rather than at the token level, thereby removing highly ``off-policy'' samples from the process of gradient calculation. This approach aligns more closely with the goals of sequence-level reward assignment and optimization. Additionally, we employ length normalization in $s_{i}(\theta)$, which serves to stabilize its variance and to maintain $s_{i}(\theta)$ within a consistent numerical range. Without such normalization, even small changes in token likelihoods could cause significant swings in the sequence-level importance ratio, necessitating different clipping ranges for responses of varying lengths. It should also be highlighted that the typical clipping ranges for GSPO, as opposed to earlier methods such as GRPO, usually vary greatly in scale because importance ratios are defined differently in each approach.

\subsection{Gradient Analysis}
\label{subsec:gradient}

We can derive the gradient of the GSPO objective as follows (clipping is omitted for brevity):

\begin{align}
\nabla_{\theta} \mathcal{J}_\text{GSPO} (\theta)
=&\ 
\nabla_{\theta} \mathbb{E}_{ x \sim \mathcal{D},\, \{y_i\}_{i=1}^G \sim \pi_{\theta_\text{old}}( \cdot | x) }
\left[ 
\frac{1}{G} \sum_{i=1}^{G} s_{i}(\theta) \widehat{A}_{i}
\right] \\
= &\ 
\mathbb{E}_{ x \sim \mathcal{D},\, \{y_i\}_{i=1}^G \sim \pi_{\theta_\text{old}}( \cdot | x) }
\left[ 
\frac{1}{G} \sum_{i=1}^{G}
s_{i}(\theta) \widehat{A}_{i}
\cdot \nabla_{\theta} \log s_{i}(\theta)
\right] \\
=&\ 
\mathbb{E}_{ x \sim \mathcal{D},\, \{y_i\}_{i=1}^G \sim \pi_{\theta_\text{old}}( \cdot | x) }
\left[ 
\frac{1}{G} \sum_{i=1}^{G} 
\left( \frac{ \pi_{\theta} (y_i | x) }{ \pi_{\theta_\text{old}} (y_i | x)} \right)^{\frac{1}{|y_i|}} \widehat{A}_{i} 
\cdot \frac{1}{|y_i|} \sum_{t=1}^{|y_i|} \nabla_{\theta} \log \pi_{\theta} (y_{i,t} | x, y_{i,<t}) 
\right].
\label{equ:gspo_grad}
\end{align}

For comparison, the gradient of the GRPO objective is as follows (note that $\widehat{A}_{i,t} = \widehat{A}_{i}$):

\begin{align}
\nabla_{\theta} \mathcal{J}_\text{GRPO} (\theta) 
=&\ 
\nabla_{\theta} \mathbb{E}_{ x \sim \mathcal{D},\, \{y_i\}_{i=1}^G \sim \pi_{\theta_\text{old}}( \cdot | x) }
\left[ \frac{1}{G} \sum_{i=1}^{G} \frac{1}{|y_i|} \sum_{t=1}^{|y_i|} 
w_{i,t}(\theta) \widehat{A}_{i,t}
\right] \\
=&\
\mathbb{E}_{ x \sim \mathcal{D},\, \{y_i\}_{i=1}^G \sim \pi_{\theta_\text{old}}( \cdot | x) }
\left[ \frac{1}{G} \sum_{i=1}^{G} \widehat{A}_{i} 
\cdot \frac{1}{|y_i|} \sum_{t=1}^{|y_i|} 
\frac{ \pi_{\theta} (y_{i,t} | x, y_{i,<t}) }{ \pi_{\theta_\text{old}} (y_{i,t} | x,y_{i,<t})} 
\nabla_{\theta} \log \pi_{\theta} (y_{i,t} | x, y_{i,<t})  
\right].
\end{align}

The primary difference between GSPO and GRPO centers on \textit{the manner in which they assign weights to the gradients of token log likelihoods}. GRPO applies a weighting scheme to each token based on its corresponding ``importance weight'' $\frac{ \pi_{\theta} (y_{i,t} | x, y_{i,<t}) }{ \pi_{\theta_\text{old}} (y_{i,t} | x,y_{i,<t})}$. These weights are not uniform; depending on the value of $\widehat{A}_i$, they may fall within $(0, 1+\varepsilon]$ when $\widehat{A}_i > 0$, or within $[1-\varepsilon, +\infty)$ for $\widehat{A}_i < 0$. Such disparities in weighting are substantial and, as training proceeds, their cumulative effect can introduce erratic behavior. Conversely, GSPO applies a uniform weighting to every token in a response, thus removing the instability associated with GRPO’s uneven weighting approach.

\subsection{GSPO-token: A Token-level Objective Variant}

In scenarios like multi-turn RL, we may desire a finer-grained advantage adjustment than the sequence level. To this end, we introduce a token-level objective variant of GSPO, namely \textbf{GSPO-token}, to allow token-wise advantage customization:

\begin{align}
\mathcal{J}_\text{GSPO-token}(\theta)
=
\mathbb{E}_{ x \sim \mathcal{D},\, \{y_i\}_{i=1}^G \sim \pi_{\theta_\text{old}}( \cdot | x) }
\left[ 
\frac{1}{G} \sum_{i=1}^{G} \frac{1}{|y_i|} \sum_{t=1}^{|y_i|} 
\min \left( s_{i,t}(\theta) \widehat{A}_{i,t},  \, \mathrm{clip} \left( s_{i,t}(\theta), 1 - {\varepsilon}, 1 + {\varepsilon} \right) \widehat{A}_{i,t} \right) 
\right],
\label{equ:gspo-token}
\end{align}

where

\begin{align}
s_{i,t}(\theta) 
= \mathrm{sg} \left[ s_{i}(\theta) \right]  \cdot \frac{ \pi_{\theta} (y_{i,t} | x, y_{i,<t}) }{ \mathrm{sg} \left[ \pi_{\theta} (y_{i,t} | x, y_{i,<t}) \right] },
\end{align}

and $\mathrm{sg}[\cdot]$ denotes only taking the numerical value but stopping the gradient, corresponding to the \texttt{detach} operation in PyTorch. The gradient of GSPO-token can be derived as:

\begin{align}
\nabla_{\theta} \mathcal{J}_\text{GSPO-token}(\theta)
=&\ 
\nabla_{\theta} \mathbb{E}_{ x \sim \mathcal{D},\, \{y_i\}_{i=1}^G \sim \pi_{\theta_\text{old}}( \cdot | x) }
\left[ 
\frac{1}{G} \sum_{i=1}^{G}
\frac{1}{|y_i|} \sum_{t=1}^{|y_i|} 
s_{i,t}(\theta) \widehat{A}_{i,t}
\right] \\
=&\ 
\mathbb{E}_{ x \sim \mathcal{D},\, \{y_i\}_{i=1}^G \sim \pi_{\theta_\text{old}}( \cdot | x) }
\left[ 
\frac{1}{G} \sum_{i=1}^{G} s_i (\theta) \cdot \frac{1}{|y_i|} \sum_{t=1}^{|y_i|} 
\widehat{A}_{i,t} \frac{ \nabla_{\theta} \pi_{\theta} (y_{i,t} | x, y_{i,<t}) }{ \pi_{\theta} (y_{i,t} | x, y_{i,<t}) }
\right] \\
=&\ 
\mathbb{E}_{ x \sim \mathcal{D},\, \{y_i\}_{i=1}^G \sim \pi_{\theta_\text{old}}( \cdot | x) }
\left[ 
\frac{1}{G} \sum_{i=1}^{G}  \left( \frac{ \pi_{\theta} (y_i | x) }{ \pi_{\theta_\text{old}} (y_i | x)} \right)^{\frac{1}{|y_i|}}
\cdot \frac{1}{|y_i|} \sum_{t=1}^{|y_i|}
\widehat{A}_{i,t} \nabla_{\theta} \log \pi_{\theta} (y_{i,t} | x, y_{i,<t})  
\right].
\label{equ:gspo-token_grad}
\end{align}

It should be observed that $\frac{ \pi_{\theta} (y_{i,t} | x, y_{i,<t}) }{ \mathrm{sg} \left[ \pi_{\theta} (y_{i,t} | x, y_{i,<t}) \right] }$ evaluates to 1, implying that $s_{i,t}(\theta)$ is equivalent in value to $s_{i}(\theta)$. When we examine Equation~\eqref{equ:gspo} alongside Equation~\eqref{equ:gspo-token}, as well as Equation~\eqref{equ:gspo_grad} with Equation~\eqref{equ:gspo-token_grad}, we find that GSPO-token yields identical numerical results for the optimization objective, clipping rule, and gradient formulation as GSPO, provided that the advantages assigned to each token in $y_i$ are uniform (specifically, $\widehat{A}_{i,t} = \widehat{A}_{i}$). However, GSPO-token offers greater versatility by permitting the assignment of different advantages to individual tokens.

\section{Experiments and Discussion}

\subsection{Empirical Results}

We conduct experiments utilizing a cold-start model fine-tuned from Qwen3-30B-A3B-Base, presenting both the training reward trajectories and model performance plots across the AIME'24 (computed as average Pass@1 over 32 samples), LiveCodeBench (spanning 202410-202502, with average Pass@1 across 8 samples), and CodeForces (assessed via Elo Rating) evaluation sets. For reinforcement learning (RL) training, we divide each rollout batch into four mini-batches to perform gradient updates. In the context of GSPO, the left and right clipping thresholds within Equation~\eqref{equ:gspo} are specified as 3e-4 and 4e-4, respectively. Our baseline for comparison is GRPO, where the clipping boundaries in Equation~\eqref{equ:grpo} are adjusted to 0.2 and 0.27, following systematic tuning to enable equitable benchmarking. It is important to highlight that GRPO requires the adoption of the Routing Replay training procedure to ensure typical convergence for MoE RL, as further elaborated in \S~\ref{subsec:routing_replay}. In contrast, \textit{GSPO eliminates the reliance on this approach}.

As depicted in Figure~\ref{fig:results}, training with GSPO exhibits consistent stability across the process. Our findings indicate that \textbf{GSPO facilitates ongoing enhancements in performance by increasing training compute, systematically refreshing the query set, and expanding generation length}. In addition, GSPO attains higher training efficiency compared to GRPO, evidenced by improved training accuracy and benchmark results for equivalent compute and query consumption. Ultimately, GSPO has been effectively utilized in RL training for recent Qwen3 models, providing compelling support for its capability to unlock the full potential of RL scaling within large language models.

\subsection{Curious Observation on Clipping Fractions}

One notable difference between GSPO and GRPO lies in how clipping is applied: GSPO clips entire sequences, whereas GRPO clips individual tokens. As illustrated in Figure~\ref{fig:clipping}, the proportion of clipped tokens for GSPO differs by two orders of magnitude from that of GRPO, and varying the clipping ranges does not mitigate this significant discrepancy. Yet, despite GSPO discarding substantially more tokens—resulting in fewer tokens being available for model optimization or gradient calculation—it consistently demonstrates superior training efficiency over GRPO. This surprising outcome, where greater token clipping translates into improved training performance, suggests that the gradient estimates produced at the token level by GRPO are fundamentally less reliable and are characterized by elevated noise, undermining their utility for sample exploitation. GSPO’s focus on entire sequences, on the other hand, produces a learning signal that is both more stable and more conducive to effective training.

\subsection{Benefit of GSPO for MoE Training}
\label{subsec:routing_replay}

\paragraph{Background}
The reinforcement learning (RL) training process for dense models contrasts sharply with that for Mixture-of-Experts (MoE) architectures, which activate only a subset of parameters at each step and consequently face distinct stability issues. Specifically, in our investigations using the GRPO algorithm, we observed that the inherent \textit{expert-activation volatility} within MoE frameworks significantly impedes the convergence of RL optimization. To elaborate, post gradient updates, the selection of experts activated for a given response can fluctuate widely. For instance, when employing the 48-layer Qwen3-30B-A3B-Base model, analysis across successive RL gradient steps reveals that, for the same sample, approximately 10\% of the experts chosen by the updated policy $\pi_{\theta}$ differ from those selected by the previous policy $\pi_{\theta_\text{old}}$. This volatility intensifies with increased model depth and leads to severe oscillations in the token-level importance weights $w_{i,t}(\theta) = \frac{ \pi_{\theta} (y_{i,t} | x, y_{i,<t}) }{ \pi_{\theta_\text{old}} (y_{i,t} | x,y_{i,<t})}$, ultimately undermining the reliability of these ratios, as examined in \S~\ref{sec:motivation} and~\ref{subsec:gradient}. These unstable conditions, in turn, obstruct effective RL convergence.

\paragraph{Our Previous Approach} In addressing this issue, our earlier solution involved leveraging the \textbf{Routing Replay} training paradigm. This method entails recording the experts activated under $\pi_{\theta_\text{old}}$, and then ``replaying'' the identical routing selections during the computation of importance ratios in $\pi_{\theta}$, specifically $w_{i,t}(\theta) = \frac{ \pi_{\theta} (y_{i,t} | x, y_{i,<t}) }{ \pi_{\theta_\text{old}} (y_{i,t} | x,y_{i,<t})}$. By ensuring both $\pi_{\theta}$ and $\pi_{\theta_\text{old}}$ utilize the same experts for each token $y_{i,t}$, we maintain identical network activations, thereby recovering the consistency and stability of token-level importance ratios, which further guarantees consistent optimization of the active subnetwork throughout gradient-based training steps. As illustrated in Figure~\ref{fig:routing_replay}, Routing Replay plays a pivotal role in achieving reliable convergence of GRPO when training MoE architectures.

\paragraph{Benefit of GSPO} While Routing Replay permits the GRPO training of MoE models to achieve convergence, its approach of maintaining previous routing modes introduces extra demands on memory and communication resources and can restrict the MoE model's effective capacity. In comparison, as illustrated in Figure~\ref{fig:results}, GSPO dispenses with reliance on Routing Replay, supports direct computation of importance ratios $s_i(\theta)$ using standard methods, and ensures stable and proper convergence during optimization. The principal advantage arises from GSPO’s exclusive emphasis on sequence likelihood (specifically, $\pi_{\theta} (y_i | x)$) rather than focusing on per-token likelihood (such as $\pi_{\theta} (y_{i,t} | x, y_{i,<t})$). Because the MoE model consistently retains its ability to perform language modeling, the sequence likelihood remains relatively stable. Ultimately, GSPO addresses the problem of expert-activation instability in MoE models at its core, eliminating the necessity for cumbersome mechanisms like Routing Replay. As a result, this approach not only streamlines and secures the training process but also ensures the model operates at its full capacity without imposed limitations.

\subsection{Benefit of GSPO for RL Infrastructure}

Due to the differences in numerical precision between various training engines (such as Megatron) and inference engines (like SGLang and vLLM), practitioners conventionally recalculate the likelihoods of generated responses using the training engine for the previous policy $\pi_{\theta_\text{old}}$. In contrast, GSPO optimizes based on sequence-level likelihoods rather than token-level ones; as a result, sequence-level likelihoods are generally less sensitive to discrepancies in precision. This design allows GSPO to leverage the likelihoods computed by the inference engine directly during optimization, eliminating the need for additional recomputation by the training engine. Such a feature proves advantageous, particularly in use cases such as partial rollout, multi-turn RL, or frameworks where training and inference are decoupled.

\section{Conclusion}

We introduce Group Sequence Policy Optimization (GSPO), a novel reinforcement learning algorithm tailored for training large language models. Adhering to the core tenets of importance sampling, GSPO computes importance ratios using sequence likelihood and applies clipping, reward assignment, and optimization at the sequence level. Our empirical results show that GSPO offers significantly enhanced stability, training efficiency, and overall performance compared to GRPO, and proves particularly effective for large-scale RL training of MoE models. This method has served as the basis for remarkable gains in the latest Qwen3 models. Leveraging GSPO as a scalable algorithmic framework, we plan to further extend reinforcement learning and anticipate substantial progress in the development of more intelligent systems.

\end{document}